% !TEX root = ../main_zh.tex
\chapter{价格、实际产出与通胀指数}
\label{ch:prices}

第~\ref{ch:accounts}~章的国民经济核算用一个数字概括了一个经济体一年内生产的全部产品的价值。可这个数字是用货币来衡量的，而货币是一把会滑动的尺子：GDP 上升 10\% 既可能意味着经济体多生产了 10\% 的产品，也可能意味着产量分毫未变、只是价格涨了 10\%，或是介于二者之间的任何组合。要判断一个国家究竟有没有增长，要把它和自己的过去相比，或是和另一个国家相比，我们都得先把支出变化中反映“产出更多”的那一部分，同反映“价格更高”的那一部分区分开来。把这两者拆开，正是本章要做的事。

我们分两步走。第一步，把名义 GDP 拆成一个实际数量和一个价格水平，由此得到 GDP 平减指数（GDP deflator），并把名义增长干净利落地分解为实际增长加上通货膨胀。有了这个工具，我们终于可以对 GDP 做跨期和跨国的比较；其间我们会停下来讨论几个实务上的报告惯例——同比与环比、以及经季节调整的折年率——正是它们使公布的增长数字之间具有可比性。第二步，我们梳理经济学家实际查阅的那一族价格指数：消费者价格指数（CPI）及其“核心”变体、生产者价格指数（PPI），以及个人消费支出（PCE）指数。一个反复出现的主题是：并不存在唯一的、那个“总”价格水平；每个指数回答的是不同的问题，选错了指数，就会悄无声息地扭曲答案。最后我们把 CPI 和 GDP 平减指数直接放在一起比较，因为二者之间的差异——固定权重对变动权重、进口品对国产品——恰恰是会出现在真实政策辩论中的那类测量细节。

关于本章范围的一点说明。本章谈的是\emph{度量}价格与通胀，而非\emph{解释}它们。整体价格水平为什么会变动、通胀给社会带来什么代价、一个靠印钞的政府所征收的“通胀税”，以及把货币增长同通胀联系起来的货币数量论，这些都留到第~\ref{ch:quantity}~章再讲。这里我们只打造尺子；至于被度量的东西为什么会动，其背后的理论留待之后。

\section{名义 GDP 与实际 GDP}

名义 GDP 从一年到下一年的变化里，混杂着两件本应分开的事：数量可能变了，价格也可能变了。要把数量单独分离出来，我们就把产出按某个选定的\emph{基期}的价格固定下来重新计算。这正是名义 GDP 与实际 GDP 之分背后唯一的那个想法。

\defn{名义 GDP 与实际 GDP}{
\emph{名义 GDP}（nominal GDP）按每种产品在其生产当年的现行价格来计价，
\[
\text{名义 GDP}_t = \sum_i p_{i,t}\, q_{i,t},
\]
其中 $p_{i,t}$ 与 $q_{i,t}$ 是产品 $i$ 在第 $t$ 年的价格与数量。\emph{实际 GDP}（real GDP）则把同样的这些数量按某个基期 $0$ 的固定价格来计价，
\[
\text{实际 GDP}_t = \sum_i p_{i,0}\, q_{i,t}.
\]
因为实际 GDP 把价格固定住了，它的每一处变化反映的都是实物产出的变化，而非价格的变化。
}

\rmkb[中国的数据口径]{中国国家统计局公布这两套序列时，用的名称翻译过来分别是“现价 GDP”，也就是名义 GDP；以及“不变价 GDP”，也就是按某个声明的基期来度量的实际 GDP。通常所说的“GDP 增速”和公布的“GDP 指数”指的都是\emph{实际} GDP 的增长。读到任何这类序列时，要问的第一个问题永远是：以哪一年为基期，以及这是一个现价还是不变价的数字？}

\subsection{GDP 平减指数}

如果说实际 GDP 抓住的是数量，那么拿名义 GDP 同实际 GDP 一比，剩下的那部分就必然抓住的是价格。这个余项就是整个经济体涵盖面最广的价格指数。

\defn{GDP 平减指数}{
\emph{GDP 平减指数}是名义 GDP 与实际 GDP 之比，按惯例乘以 $100$ 来标度，
\[
\text{GDP 平减指数}_t = \frac{\text{名义 GDP}_t}{\text{实际 GDP}_t}\times 100
= \frac{\sum_i p_{i,t}\, q_{i,t}}{\sum_i p_{i,0}\, q_{i,t}}\times 100 .
\]
它是一个涵盖进入 GDP 的\emph{所有}产品的价格指数，并以\emph{当年}的数量为权重。
}

在基期，名义 GDP 与实际 GDP 重合，平减指数恰好等于 $100$。只要价格自基期以来平均而言是上升的——在任何存在正通胀的经济体里都属常态——分子就大于分母，平减指数便在 $100$ 之上；而在基期之前，它会落在 $100$ 以下。一个好记的口诀是：平减指数在多数时候都大于 $100$，原因恰恰在于基期通常落在过去。

\subsection{名义增长的分解}

现在我们可以把“名义增长就是实际增长加上通货膨胀”这句含糊的话讲精确了。记 $g^{N}$ 为名义 GDP 的增长率、$g^{R}$ 为实际 GDP 的增长率、$\pi$ 为 GDP 平减指数的增长率（即平减指数通胀率）。由于名义 GDP 是实际 GDP 与平减指数的乘积，它的毛增长因子就是另外两者的乘积：

\thm{增长分解}{
毛增长因子满足
\[
1 + g^{N} = \pr{1 + g^{R}}\pr{1 + \pi}.
\]
展开后丢掉二阶交叉项 $g^{R}\pi$——当两个增长率都很小时这一项也很小——就得到便于使用的近似式
\[
g^{N} \approx g^{R} + \pi .
\]
}

\pf{
按定义有 $\text{名义 GDP} = \text{实际 GDP}\times(\text{平减指数}/100)$，于是在相邻两年间取比值，
\[
1 + g^{N}
= \frac{\text{名义 GDP}_{t+1}}{\text{名义 GDP}_t}
= \frac{\text{实际 GDP}_{t+1}}{\text{实际 GDP}_t}\cdot\frac{\text{平减指数}_{t+1}}{\text{平减指数}_t}
= \pr{1 + g^{R}}\pr{1 + \pi}.
\]
乘开后 $1 + g^{N} = 1 + g^{R} + \pi + g^{R}\pi$，故 $g^{N} = g^{R} + \pi + g^{R}\pi$。对于量级为百分之几的增长率，乘积 $g^{R}\pi$ 的量级仅在百分之零点几，故予以舍去。
}

这个分解写成乘积形式时是精确的，写成加法形式时则只是近似；两者的差距就是那个交叉项，它只在增长率或通胀率很大时才要紧。这与支配利息的复利逻辑是同一回事：基期一旦确定，在正常年份里，名义 GDP 会越来越偏离实际 GDP，平减指数也越来越大，因为每一年的价格上涨都在前一年的基础上复利累积。

\section{GDP 的跨期与跨国比较}

我们之所以费这番功夫把价格从 GDP 里剥离出来，是因为我们想要\emph{比较}：今年比去年、一国比另一国。正是实际 GDP 让这些比较变得有意义。

\subsection{跨期比较与跨国比较}

要把一个经济体同自己的过去相比，就必须把价格因素去掉，而这正是实际 GDP 所做的：实际 GDP 的上升是产品和服务数量上的真实增加，而不是通胀造成的假象。要做跨国比较，我们还需额外把数字换算成同一种货币（并且，做福利比较时，往往还要通过购买力平价来修正价格水平的差异，这一精细化处理在此暂且不论）。于是 GDP 既为追踪长期增长、又为在某一时点上给经济体排序，提供了一个共同的基准。

即便在使用这把尺子时，也值得把它的局限放在眼前。GDP 记录的是生产的市场价值；它对这份产出如何分配、对财富（而非收入）、对环境退化、对腐败，以及对人们是否因这些被它计入的活动而更快乐，统统是缄默的。这些都不会让 GDP 变得无用——跨国家、跨数十年来看，从预期寿命到识字率，GDP 同我们真正在乎的那些东西都稳健地正相关。正确的结论不是 GDP 可以被抛弃、或被某个更好的单一数字取代，而是人类福祉本就是多维的，最好同时参照若干个指标来读，而 GDP 是其中信息含量最高的指标之一。

\subsection{同比、环比与折年率}

高频数据带来了一个年度数字会掩盖的报告难题。某一季度的产出，既可以同\emph{去年同}季度相比（同比增长），也可以同\emph{紧邻的上一}季度相比（环比增长），而这两者讲述的是不同的故事。

\begin{itemize}
    \item \textbf{同比}增长拿相同季度同相同季度相比，因而会自动抵消掉规律性的季节模式——比如说，零售额总会在某个节日前后激增。它的短处是慢：因为它把中间整整一年的变化都平均了进去，所以在揭示经济的转折点上可能反应迟缓。
    \item \textbf{环比}增长则及时，变化一发生就被它记录下来；但也正因如此，它被季节性所污染。要让它可用，就必须做\emph{季节调整}，即在计算增长率之前，先用统计方法把年内规律性的季节模式剔除掉。
\end{itemize}

\rmkb[本应吻合、却未必吻合的两套序列]{这两种口径并非彼此独立——一套自洽的季度水平数据，应当能让同比与环比增长率相互对得上——可实务中，中国公布的同比和环比 GDP 数字有时却无法相互印证，这提醒我们：季节调整这一步本身就是一种估计，对高频宏观数据应当抱着一些容忍度去读。}

有若干经济体——美国、日本和欧元区都在其列——会报告一个\emph{经季节调整的折年率}（seasonally adjusted annualized rate，SAAR）。其想法是，取当季经季节调整后的环比增长，然后追问：\emph{假如}经济按这个季度速度连续增长四个季度，会得到怎样的年增长率？写成公式，若 $g$ 为环比增长率，则折年率为
\[
\text{SAAR} = \pr{1+g}^{4} - 1 .
\]
折年率既及时又直观——它把一个季度的读数用人们熟悉的年率单位讲出来——但由于它把单个季度的增长复利了四次，它也把那个季度的噪声放大了，所以 SAAR 比同比增长波动得厉害得多。

\rmkb[被熨平的数据与“伤疤效应”]{统计机构可能会刻意“削峰填谷”，让公布出来的序列看上去比真实经济更平滑。熨平\emph{产出}序列是一回事——GDP 还吸收得了——但底层那个砸向\emph{收入}的冲击却没那么容易被糊弄过去：一次真正的衰退会永久性地削减受影响者的收入，而支出下跌得还要更多，一方面是因为收入降了，另一方面是因为家庭会以预防性储蓄来应对冲击，从而持久地改变其消费与储蓄习惯。一次急剧的收缩在收缩本身过去很久之后，仍留在经济活动水平上的那道持久创伤，就叫做\emph{伤疤效应}（scar effect）。}

\section{实务中的价格指数}

GDP 平减指数涵盖面广，但它公布得不勤、又有滞后，而且它也不是观察生活成本、或观察对货币政策最为相关的那些压力时被盯得最紧的指数。实务中，经济学家会查阅一小族价格指数，每一个都是为不同的目的而构造的。我们逐个来看。

\subsection{消费者价格指数}

\defn{消费者价格指数（CPI）}{
\emph{消费者价格指数}度量的是一\emph{篮固定的}、能代表典型家庭消费的产品与服务的成本。每一类消费被赋予一个反映其在代表性家庭预算中占比的权重，指数随时间追踪这一篮固定的、加权的产品的成本。一言以蔽之，它是对生活成本的一种衡量。
}

CPI 的构造就是为了回答一个问题：本月再去买典型家庭从前买的那同一篮消费品，要多花（或少花）多少？由于这一篮商品及其权重被固定住了——在中国，它们大约每五年才修订一次，而且各细类的权重并不对外公布，不过可以从其他公布的数字间接倒推出来——这个指数就把价格的变化同人们所买东西的变化分离开了。

\rmkb[篮子里装了什么——一幅中国快照]{食品在中国 CPI 里占很大的权重；单是猪肉，历史上就一直是权重最高的单一商品，而由于养猪遵循一个政策很难驯服的生产周期，猪肉价格大幅波动，把指数也一并拖着走。范围更广的食品烟酒类大约占指数的四分之一。这里有一条一般性的规律：食品消费占比\emph{越低}，生活水准就\emph{越高}——这就是恩格尔定律，即越富裕的家庭，花在食品上的收入比例越小，花在一切能提升生活品质的东西上的比例越大。居住主要通过\emph{租金}进入 CPI：因为中国自有住房比例高，买房被当作投资而非消费，所以房\emph{价}的涨跌并不直接进入 CPI，于是房地产价格的迅速攀升或崩塌可能根本不会在指数里显现。出于同样的道理，股票和债券价格——它们是资产，不是消费品——也被排除在外。}

一个标准惯例把 CPI 通胀高于 $2\%$ 看作真正的通货膨胀，把低于 $2\%$ 看作通胀放缓、渐渐过渡到通货紧缩。CPI 在实务上最大的优点是频率：它每月公布，这使它在实时分析上远比按季公布、又有滞后的平减指数有用。

\rmkb[CPI 并非全貌]{只靠 CPI 来讨论经济，得到的是一幅不完整的图景。直到 2000 年代中期，美国 CPI 一直波澜不惊，而房价和次级抵押贷款市场却剧烈波动；这个指数根本没看到那场最终在 2008 年酿成危机的金融风险的累积。事后吸取的教训是：必须把消费价格同资产价格——房地产和金融资产——放在一起看，决策者必须留意一组更广的宏观指标、留意系统性风险，这就是如今所称的\emph{宏观审慎}（macroprudential）政策立场。}

\subsection{成本推动型与需求拉动型通货膨胀}

价格为什么上涨，与它涨得有多快同等要紧，因为压力的来源决定了政策能否对它有所作为。把两条渠道分开来看是有益的。

\defn{成本推动型与需求拉动型通货膨胀}{
\emph{成本推动型通货膨胀}（cost-push inflation）源自\emph{供给}侧：投入品成本——能源、食品、原材料——的上升，即便在需求不变时也把价格往上推。\emph{需求拉动型通货膨胀}（demand-pull inflation）源自\emph{需求}侧：更旺盛的支出，把价格拉着，去对抗只能缓慢调整的供给。
}

这两者在数据里有各自的特征印记。食品和能源价格在很大程度上由供给状况驱动——收成、天气、石油市场——所以能源价格的急剧上升正是成本推动型通胀的教科书案例。相比之下，制造业产出是在相对稳定的供给条件下生产的，所以制成品价格的变动通常被归因于需求的变化，这是需求拉动型通胀的教科书案例。这一区分不只是描述性的：我们将会看到，它告诉我们货币政策有望应对的是哪一种通胀。

\subsection{核心 CPI 与货币政策的逻辑}

\defn{核心 CPI}{
\emph{核心 CPI}（core CPI）是把\emph{食品和能源}成分从篮子里剔除后重新计算的消费者价格指数。通过剔除那些最易波动、由供给驱动的价格，它比整体 CPI 更平滑地追踪价格的底层趋势。
}

核心 CPI 的动机，恰恰就在成本推动／需求拉动之分。食品和能源是最易受供给冲击影响的价格——即成本推动那一部分——它们也是波动最大的。把它们剔除后，留下的指数波动更小，而且就其构造而言，主要反映的是需求侧的压力，而非转瞬即逝的供给冲击。

这正是一家中央银行所需要看到的，而其缘由是结构性的。

\state{中央银行为何盯住核心通胀}{
货币政策通过总\emph{需求}起作用：宽松或收紧，改变的是经济体想要花掉多少。它\emph{不}作用于供给侧——它没法让石油变便宜、也没法让收成变大。因此，中央银行应当对它实际能够影响的那部分通胀——也就是需求拉动型通胀——做出反应。由于核心 CPI 把由供给驱动的、成本推动的那一部分剔除掉了，作为货币政策的向导，它比整体 CPI 更可靠。
}

这条原则的实际威力，在它被无视时看得最清楚；接下来两则注记记录了几个通胀来源被误诊的事例。

\rmkb[中国 2008：朝着供给冲击去收紧]{2008 年初一个异常寒冷的冬天，让中国南方大量生猪死亡；猪肉价格飙升，单月 CPI 读数一度达到 $8\%$ 上下。在那年春天的全国人民代表大会上，领导层提出要收紧基础货币的供给。但这场通胀压倒性地是成本推动型的——是一次砸向猪肉的供给侧冲击，底层需求几乎没变——所以收紧货币是用错了工具，收效甚微。更糟的是，美国金融危机当时正开始波及中国，而出口对中国经济又如此重要，审慎之举本应是准备\emph{宽松}。2008 年下半年增长明显放缓，下行压力加大，秋天的党代会便顺势提出要放松。短短一年之内立场来了个彻底反转，暴露出当时的政策框架是多么缺乏前瞻性。}

\rmkb[中国 2018 与猪肉难题]{一句广为流传的俏皮话道出了 2018 年猪肉价格的飙升：“加上猪肉，全是通胀；去掉猪肉，全是通缩。”货币政策碰不到供给侧，可面对一个很高的\emph{总体}通胀数字——以及它所引发的公众与市场反应——中央银行可能感到不得不收紧。危险在于，朝着供给驱动的通胀去收紧，会制造出\emph{滞胀}：通胀既是成本推动型的，便治不好，而经济却被进一步往下压。}

\rmkb[美国的供给侧通胀]{同一时期美国的一段经历讲出了同样的道理。那一时期的通胀是由能源短缺、劳动力供应紧张和物流系统不畅所驱动的——全都是供给侧的问题。比方说，劳动力供应之所以紧张，部分缘于优厚的福利制度压低了劳动参与率，而非需求过热。需求侧的货币政策同一个供给侧的问题严重不匹配，与其说能解决通胀，不如说有加剧通胀之虞。}

\subsection{PCE 指数}

\defn{个人消费支出（PCE）指数}{
\emph{个人消费支出}价格指数是美国联邦储备体系偏好的消费通胀度量。它的构造与 CPI 大体相同、走势也与之大体一致，但由美联储单独计算，而美联储在其政策权衡中对通胀给予极高的权重。
}

美联储宁愿自行维护一套消费价格度量，而不直接采用公布的 CPI，这件事本身就表明：一家中央银行在价格指数的选择上有多么认真。PCE 与 CPI 彼此走势贴得很近，但在覆盖范围和权重上有所不同；二者之间的差距在正常时期很小，却真实存在，而一家盯住通胀的中央银行，会在意自己手里拿的是哪一把尺子。

\subsection{生产者价格指数}

\defn{生产者价格指数（PPI）}{
\emph{生产者价格指数}度量的是\emph{出厂}环节的价格——生产者就其产出所收到的价格，而非消费者所支付的零售价格。出厂价格大致是成本加利润加税金，它也是商人为了转售而进货时所支付的批发价格。
}

CPI 和 GDP 平减指数看的都是零售价格，即买家在最终销售环节面对的价格；PPI 则往上游看一步，看的是生产者收取的价格。这使它成为一个不同而又互补的工具。PPI 不含农产品，与需求紧密相连、又与企业利润直接挂钩，因而对经济周期相当敏感，这让它对企业的投资决策颇具信息含量。

它的关键分界线是零。当 PPI 为正时，产出价格在上涨、生产者的利润空间在扩大；当它转负时，生产者利润被挤压，生产与投资的意愿随之收缩。这甚至会\emph{自我强化}：下跌的生产者价格压低利润和投资，进而削弱经济活动、把价格进一步推低，在短期内把经济拖入一个恶性循环——或者，把符号反过来，拖入一个良性循环。

\rmkb[一次通缩式的读法]{在一个具有代表性的月份——2022 年 10 月——中国 CPI 略高于 $1\%$，而 PPI 已经转负。消费价格通胀疲软、生产者价格则干脆下跌，消费需求和投资需求都很弱，经济正滑向通货紧缩。把这两个指数放在一起读，对这场放缓的诊断，比单看其中任何一个都要清晰。}

\section{CPI 与 GDP 平减指数之比较}

现在我们手上有了两个涵盖面广的价格指数——CPI 与 GDP 平减指数——值得把它们仔细比较一番，因为二者在三个要紧的方面有所不同，也因为这一比较揭示出 CPI 身上一个系统性的偏误。

\paragraph{固定权重与替代偏误。} CPI 建立在一篮\emph{固定}的商品之上：它问的是，年复一年地买相同的数量，要花多少钱，不管相对价格怎么变。GDP 平减指数则相反，它以\emph{当年}的数量为权重来给商品加权，而这些数量会随相对价格的变化而调整。这个差异是有方向的。当某种商品变得相对更贵时，家庭会从它身上替代开来、转向更便宜的替代品，所以维持给定生活水准的真实成本，其上升幅度小于那一篮旧的、固定的商品的成本上升幅度。由于把数量固定住了，CPI 无视了这种替代，因而\emph{高估}了生活成本的上升——它患有\emph{替代偏误}（substitution bias）。而平减指数通过当年权重把替代嵌了进去，便不存在此问题。出于同样的效果，CPI 对新产品的问世、以及既有产品质量的提升都反应迟缓，而这两者都在不抬高那一篮固定商品成本的情况下提升了福利。

\paragraph{进口品。} 两个指数在来源范围上也不同。GDP 平减指数只覆盖\emph{国内生产}的商品——GDP 按定义就把进口排除在外。CPI 覆盖的则是消费者实际购买的东西，其中包含进口品。比方说，进口石油价格的跳涨会直接抬高 CPI，但它进入 GDP 平减指数，仅限于它影响到国产产出价格的那部分。

\paragraph{商品的国内范围。} 最后，两个指数覆盖经济中不同的切片。CPI 只追踪消费品。GDP 平减指数追踪 GDP 中的一切——不仅是消费，还有投资、政府购买和净出口。消费价格是平减指数的一个重要构成，但也只是其中一个构成；平减指数的覆盖面要宽广得多。两个指数回答的是不同的问题，在任何给定的年份里，它们都可能彼此分道而行。

\state{两个指数怎么读}{
CPI 是一个固定篮子的生活成本指数，每月公布，包含进口品，并倾向于通过替代偏误、以及对新产品与改良产品的迟缓纳入而\emph{高估}通胀。GDP 平减指数是一个以当年为权重、只含国内、涵盖整个 GDP 的价格指数，它把替代纳入了进来，但公布得没那么勤、还有滞后。各自适用于不同的问题：看生活成本、要及时读数时用 CPI，要看实际 GDP 背后那个全经济范围的价格水平时用平减指数。
}

本章为我们配齐了尺子——名义量对实际量、平减指数，以及消费者价格指数和生产者价格指数——并给了我们一种实用的判断力，知道在何时哪一把才是对的那把。这些工具度量的，是整体价格水平及其变化率，也就是通货膨胀。至于那个水平\emph{为什么}会动、通胀对借款人和放款人各自做了什么、政府靠印钞攫取的那笔收入，以及把货币同价格联系起来的货币数量论，则是第~\ref{ch:quantity}~章的主题；不过我们现在要先转向核算账户只是隐约提到的另一个大总量——劳动力市场与失业。
